% Integration fragment. All labels start eg: and no special macros are required.
% Prior source: Differential_Polygamma_and_Lerch_Identities, Section 8;
% current Bernoulli_Tornheim_Digamma_Identities, Sections 3--6.
% Exact verification: verify_gap_four.py, gap_four_exact_results.json.
\section{The optimal initial prefix when the even orders differ by four}
\label{eg:section}

The current incoming reports prove the optimal initial-prefix statement
for the diagonal, for adjacent even orders, and for the fixed smaller
orders $p=2,4,6$ \cite{Canonical,Differential,Bernoulli,Polygamma}.
We prove the next fixed-gap family
\[
                    (p,q)=(2b-4,2b),\qquad b\ge3.
\]
Its additional range relative to those reports begins with $(p,q)=(8,12)$ and continues
with $(10,14),(12,16),\ldots$. The proof uses the strict log-convexity
of the positive coefficients of $2/x-\cot(x/2)$ to show that a
specified two-by-two boundary determinant never vanishes. We also
prove a structural statement about every pair of columns in the
fundamental recurrence for an arbitrary even gap. The latter does
not settle the higher-dimensional boundary determinant.

\subsection{Normalization and the determinant to be proved nonzero}

For even $2\le p\le q=2b$, put
\[
 A_p(n)=\zeta(p)-H_n^{(p)}=\zeta(p,n+1),\qquad
 D_{p,q}(s)=\sum_{n\ge0}A_p(n)A_q(n)(n+\tfrac12)^{-s}.
\]
The series initially converges for $\Re s>3-p-q$, and its continuation
is regular at every $s=-2r$. Write $D_r=D_{p,q}(-2r)$. We use the
inherited blocks $C_j$, characterized by the mutually inverse formulas
\begin{align}
 D_r&=\sum_{j=0}^r\frac{\binom{2r}{2j}}{2j+1}
       B_{2r-2j}(\tfrac12)C_j,\label{eg:forward}\\
 C_j&=\sum_{r=0}^j\binom{2j+1}{2r}2^{-2(j-r)}D_r.
       \label{eg:backward}
\end{align}
For $j\ge b$ they have the finite coordinate form
\[
 C_j=\sum_{k=0}^{b-1}a_{jk}Z_k,\qquad
 Z_k=\zeta(2k),\qquad Z_0=-\tfrac12.
\]
For $d=2j+1$ and $e=d-p-q+2k+1$, the inherited coefficients are
\begin{equation}\label{eg:entries}
 a_{jk}=\begin{cases}
 B_e\left\{\displaystyle
 \frac{\binom{d-q+1}{p-2k}}{d-q+1}
 +\frac{\binom{d-p+1}{q-2k}}{d-p+1}\right\},&e\ge0,\\[6pt]
 0,&e<0.
 \end{cases}
\end{equation}
Each binomial summand is omitted unless its lower index lies between
$2$ and its upper index. In particular a zero lower index is omitted.
Let
\begin{equation}\label{eg:M}
 M_{p,q}=(a_{b+r-1,k})_{1\le r\le b,\ 0\le k<b}.
\end{equation}
The first $q$ moments have full rank on the specified formal labels
\[
 C_0,\ldots,C_{b-1},Z_0,\ldots,Z_{b-1}
\]
if and only if $M_{p,q}$ is nonsingular. The complete first-$q$
determinant in this coordinate order is
$\det M_{p,q}/(4b-1)!!$, by the triangular transform
\eqref{eg:forward}. No arithmetic independence of the actual constants
is asserted by this formal rank statement.

\subsection{A strict comparison theorem for all even gaps}

Put $q-p=2c$, where $1\le c<b$, and define
\begin{equation}\label{eg:alpha}
 \alpha_n=\frac{2|B_{2n}|}{(2n)!}
          =4\sum_{m=1}^{\infty}(2\pi m)^{-2n},\qquad n\ge1.
\end{equation}
For $s\ge c+1$ set
\begin{equation}\label{eg:ABD}
 A_s=\frac{(2s)!}{(2s-2c)!},\qquad
 B_s=\frac{(2s-1)!}{(2s-2c-1)!},\qquad
 d_s=2(A_s+B_s),
\end{equation}
and put $A_s=0$ for $1\le s\le c$. The letters $A_s,B_s$ in
\eqref{eg:ABD} are rational recurrence coefficients, distinguished from
the harmonic tails $A_p(n)$ and Bernoulli numbers $B_{2n}$ by their
arguments and subscripts. Define
\begin{equation}\label{eg:kernel}
 K(s,u)=\frac{\alpha_{s-u}(B_s+A_u)}{d_s},
 \qquad s\ge c+1,\quad1\le u<s,
\end{equation}
and the fundamental array
\begin{align}
 R_{s,j}&=\delta_{sj},&&1\le s,j\le c,\notag\\
 R_{s,j}&=\sum_{u=1}^{s-1}K(s,u)R_{u,j},
       &&s\ge c+1,\quad1\le j\le c.
       \label{eg:fundamental}
\end{align}
All generated entries $R_{s,j}$, $s\ge c+1$, are strictly positive.
The earlier fixed-gap reduction gives
\begin{equation}\label{eg:inherited-reduction}
 \det M_{2(b-c),2b}
 =\frac{(-1)^{b(b-1)/2}}{4^bb!}
    \left(\prod_{s=c+1}^{b+c}d_s\right)
    \det(R_{b+i,j})_{1\le i,j\le c}.
\end{equation}
We give the coefficient mechanism below to check the convention and
signs directly; nonvanishing of the last determinant is the additional
issue addressed here when $c=2$.

\begin{lemma}[Strictly decreasing kernel ratios]\label{eg:kernel-ratios}
For every $s\ge c+1$, the quotient
\[
                   \frac{K(s+1,u)}{K(s,u)}
\]
is a strictly decreasing function of the integer $u$, $1\le u<s$.
\end{lemma}

\begin{proof}
The coefficients in \eqref{eg:alpha} are strictly log-convex. Indeed,
with $x_m=(2\pi m)^{-2}$,
\[
 \alpha_n\alpha_{n+2}-\alpha_{n+1}^2
 =16\sum_{m<\ell}x_m^nx_\ell^n(x_m-x_\ell)^2>0.
\]
The sums converge absolutely for $n\ge1$, and the positive support
contains more than one point. Thus $\alpha_{n+1}/\alpha_n$ strictly
increases with $n$. On the other hand,
\[
 \frac{K(s+1,u)}{K(s,u)}
 =\frac{d_s}{d_{s+1}}
   \frac{\alpha_{s+1-u}}{\alpha_{s-u}}
   \frac{B_{s+1}+A_u}{B_s+A_u}.
\]
The middle factor strictly decreases with $u$. The last factor is
nonincreasing because $A_u$ is nondecreasing, $B_{s+1}>B_s>0$,
and $(B_{s+1}+a)/(B_s+a)$ decreases as $a$ increases. Every factor
is positive, so their product strictly decreases.
\end{proof}

\begin{theorem}[Every generated minor of order two is negative]
\label{eg:two-minors}
Fix an even gap $2c$ with $c\ge2$. For any initial columns
$1\le j<k\le c$, the positive ratios
\[
                   \frac{R_{s,k}}{R_{s,j}},\qquad s\ge c+1,
\]
strictly decrease with $s$. Equivalently, for $c+1\le r<s$,
\begin{equation}\label{eg:minor-sign}
 \det\begin{pmatrix}R_{r,j}&R_{r,k}\\R_{s,j}&R_{s,k}\end{pmatrix}<0.
\end{equation}
\end{theorem}

\begin{proof}
Write $X_u=R_{u,j}$, $Y_u=R_{u,k}$. Suppose that the stated ratio
ordering has been proved through row $s\ge c+1$, and put
$z_s=Y_s/X_s$. This supposition is empty at $s=c+1$, so it includes
the induction base. For $u<s$ set $v_u=Y_u-z_sX_u$. Among the initial
rows, $v_j=-z_s<0$, $v_k=1>0$, and all other $v_u$ vanish. Among
the generated rows $c+1\le u<s$, the induction hypothesis gives
$v_u>0$. Thus the unique negative term is at index $j$, and every
positive term has larger index. The recurrence gives
\[
                     \sum_{u=1}^{s-1}K(s,u)v_u=0.
\]
Put $\beta_u=K(s+1,u)/K(s,u)$. Since $v_s=0$,
\begin{align*}
 Y_{s+1}-z_sX_{s+1}
 &=\sum_{u=1}^{s-1}K(s,u)\beta_uv_u\\
 &=\sum_{u=1}^{s-1}K(s,u)(\beta_u-\beta_j)v_u<0.
\end{align*}
The negative-index term vanishes; every nonzero remaining term is
negative by Lemma~\ref{eg:kernel-ratios}. Strictness already follows
from the term $u=k$. Hence $Y_{s+1}/X_{s+1}<z_s$, completing the
induction. Multiplying the ratio inequality by the two positive
denominators proves \eqref{eg:minor-sign}.
\end{proof}

The one exceptional initial negative index is essential in this proof.
Positivity of the recurrence entries alone would not imply the theorem.

\subsection{Exact determinant and reconstruction for gap four}

Now take $c=2$. Define $X_s=R_{s,1}$ and $Y_s=R_{s,2}$, so
\[
 (X_1,Y_1)=(1,0),\qquad (X_2,Y_2)=(0,1),
\]
and compute both columns by \eqref{eg:fundamental}. In this case
\begin{equation}\label{eg:d-gapfour}
 d_s=16(s-1)^2(2s-1)(2s-3),\qquad s\ge3.
\end{equation}
Set
\begin{equation}\label{eg:casoratian}
 E_s=X_{s+1}Y_s-X_sY_{s+1},\qquad s\ge3.
\end{equation}
Theorem~\ref{eg:two-minors} proves $E_s>0$ at every index. For example,
\[
 E_3=\frac1{14515200},\qquad
 E_4=\frac{61}{1170505728000},\qquad
 E_5=\frac{3419}{69528040243200000}.
\]
These are exact rational recurrence values, not approximations.

\begin{theorem}[Optimal prefix for every gap-four pair]\label{eg:main}
For every integer $b\ge3$,
\begin{equation}\label{eg:determinant}
 \boxed{\displaystyle
 \det M_{2b-4,2b}
 =(-1)^{b(b-1)/2+1}
   \frac{4^b((b+1)!)^2(2b+3)!!(2b+1)!!}{3b!}\,E_{b+1}\ne0.}
\end{equation}
Consequently, for $(p,q)=(2b-4,2b)$,
\begin{align}\label{eg:span}
 \operatorname{span}_{\mathbb Q}\{D_0,D_1,\ldots\}
 &=\operatorname{span}_{\mathbb Q}\{D_0,\ldots,D_{q-1}\}\notag\\
 &=\operatorname{span}_{\mathbb Q}
     \{C_0,\ldots,C_{b-1},Z_0,\ldots,Z_{b-1}\}.
\end{align}
The complete first-$q$ coefficient determinant is
\eqref{eg:determinant} divided by $(4b-1)!!$. The prefix is optimal
for rank $q$ on these $q$ formal labels.
\end{theorem}

\begin{proof}
The last determinant in \eqref{eg:inherited-reduction} is now
\[
 \det\begin{pmatrix}X_{b+1}&Y_{b+1}\\X_{b+2}&Y_{b+2}\end{pmatrix}
 =-E_{b+1}<0.
\]
The elementary product evaluation
\[
 \prod_{s=3}^{b+2}d_s
 =\frac{16^b((b+1)!)^2(2b+3)!!(2b+1)!!}{3}
\]
gives \eqref{eg:determinant}. The block decomposition and invertible
triangular transforms give \eqref{eg:span}. Fewer than $q$ rows cannot
have rank $q$ on $q$ specified formal coordinates.
\end{proof}

The first five determinants are
\[
\begin{array}{c|c}
(p,q)&\det M_{p,q}\\\hline
(2,6)&61/5760\\
(4,8)&-3419/138240\\
(6,10)&-713039/6451200\\
(8,12)&263589503/309657600\\
(10,14)&818219064067/78033715200
\end{array}
\]
The first three pairs overlap the previously settled fixed smaller
orders. The last two belong to the additional infinite range covered
by Theorem~\ref{eg:main}.

\paragraph{Checking the source reduction and the rotation.}
For an arbitrary coefficient vector $v=(v_0,\ldots,v_{b-1})$, set
$Q(t)=\sum_{k=0}^{b-1}v_k t^{2b-2k}/(2b-2k)!$ and write
$Q(ix)=R(x)=\sum_{u=1}^b r_ux^{2u}$. Thus
$r_u=(-1)^uv_{b-u}/(2u)!$. Consider the real odd expression
\[
 (R^{(2c)}-R^{(2c)}(0))\cot(x/2)+(R\cot(x/2))^{(2c)}.
\]
The coefficient at $x^{2r-1}$, with $s=r+c$, is
\begin{equation}\label{eg:coefficient-equation}
 Y_r^*=d_sr_s-\sum_{u=1}^{s-1}\alpha_{s-u}(B_s+A_u)r_u,
 \qquad r_{b+1}=\cdots=r_{b+c}=0.
\end{equation}
Indeed $2/x$ supplies $d_s$, and the remaining cotangent coefficients
are $-\alpha_n$. The constant derivative subtraction sets $A_c=0$,
exactly enforcing the omitted zero lower binomial index.
Comparing factorials directly with \eqref{eg:entries} gives the row scaling
\[
 Y_r^*=\frac{2(-1)^{r-c}}{(2r-1)!}(M_{p,q}v)_r.
\]
Here is a complete determinant check of
\eqref{eg:inherited-reduction}. Let $L$ be the lower triangular
matrix of size $b+c$ whose first $c$ rows are the corresponding
identity rows, and whose remaining entries are
\[
 L_{ss}=d_s,\qquad L_{su}=-d_sK(s,u)\ (u<s),\qquad
 L_{su}=0\ (u>s).
\]
Its inverse has entries $(L^{-1})_{s,j}=R_{s,j}$ in its first $c$
columns. The coefficient matrix $H$ of
\eqref{eg:coefficient-equation} is the minor with rows
$c+1,\ldots,b+c$ and columns $1,\ldots,b$. Jacobi's complementary
minor identity therefore gives
\[
 \det H=(-1)^{bc}\left(\prod_{s=c+1}^{b+c}d_s\right)
                     \det(R_{b+i,j})_{1\le i,j\le c}.
\]
If $T$ denotes the reversed diagonal map from $v$ to $r$ and $S$
the row scaling just displayed, then $HT=SM_{p,q}$ and
\[
 \det T=\frac{(-1)^b}{\prod_{u=1}^b(2u)!},\qquad
 \det S=\frac{2^b(-1)^{b(b+1)/2-bc}}
                       {\prod_{r=1}^b(2r-1)!}.
\]
Taking determinants yields exactly
\eqref{eg:inherited-reduction}. This checks the recurrence against
the original Bernoulli matrix, including the sign convention, without
a new analytic assumption.

For $c=2$, these equations give a particularly small reconstruction.
Given the actual blocks $C_b,\ldots,C_{2b-1}$, set
\[
 Y_r^*=\frac{2(-1)^{r-2}}{(2r-1)!}C_{b+r-1},\qquad1\le r\le b.
\]
Set $\eta_1=\eta_2=0$ and compute, for $3\le s\le b+2$,
\begin{equation}\label{eg:inhomogeneous}
 \eta_s=\frac{Y_{s-2}^*}{d_s}
          +\sum_{u=1}^{s-1}K(s,u)\eta_u.
\end{equation}
The two initial coefficients are recovered by
\begin{equation}\label{eg:reconstruction}
 \begin{pmatrix}r_1\\r_2\end{pmatrix}
 =-\begin{pmatrix}X_{b+1}&Y_{b+1}\\X_{b+2}&Y_{b+2}\end{pmatrix}^{-1}
       \begin{pmatrix}\eta_{b+1}\\\eta_{b+2}\end{pmatrix}.
\end{equation}
Then $r_s=X_sr_1+Y_sr_2+\eta_s$ and
$Z_k=(-1)^{b-k}(2b-2k)!r_{b-k}$. This uses one two-by-two inverse
with a proved nonzero determinant. Equations \eqref{eg:backward} and
\eqref{eg:reconstruction} reconstruct all the even-zeta coordinates
from the first $q$ moments; substitution into \eqref{eg:entries} and
\eqref{eg:forward} recovers every later moment.

\subsection{An infinite family of ordinary rational polygamma sums}

Write the first row of $M_{p,q}^{-1}$ as
$(\lambda_b,\ldots,\lambda_{q-1})$ and define
\begin{equation}\label{eg:ordinary-weights}
 w_r=\sum_{j=\max\{b,r\}}^{q-1}
          \lambda_j\binom{2j+1}{2r}2^{-2(j-r)},\qquad
 P_b(x)=\sum_{r=0}^{q-1}w_rx^{2r}.
\end{equation}
For $p=q-4$, put
\begin{align*}
 u_0&=\frac1{(p-1)(q-1)},\\
 u_1&=-\frac1{24}\left(\frac q{p-1}+\frac p{q-1}\right),\\
 u_2&=\frac7{5760}\left(
 \frac{q(q+1)(q+2)}{p-1}+\frac{p(p+1)(p+2)}{q-1}\right)
          +\frac{pq}{576},\\
 S_b(x)&=u_0w_{q-1}x^4
       +(u_0w_{q-2}+u_1w_{q-1})x^2
       +u_0w_{q-3}+u_1w_{q-2}+u_2w_{q-1}.
\end{align*}

\begin{corollary}[Gap-four rational evaluations]\label{eg:ordinary-family}
For every even $q\ge6$, with $p=q-4$ and $x_n=n+1/2$,
\begin{equation}\label{eg:ordinary-family-equation}
 \boxed{\displaystyle
 \sum_{n=0}^{\infty}\left\{
 \frac{P_b(x_n)}{(p-1)!(q-1)!}
   \psi^{(p-1)}(n+1)\psi^{(q-1)}(n+1)-S_b(x_n)
                  \right\}=-\frac12.}
\end{equation}
The sum converges absolutely. Its polynomial $P_b$ has exact degree
$2q-2$, because
\begin{equation}\label{eg:top-coefficient}
 w_{q-1}=\frac{4b(4b-1)E_b}{d_{b+2}E_{b+1}}>0.
\end{equation}
\end{corollary}

\begin{proof}
The inverse row gives $\sum_{j=b}^{q-1}\lambda_jC_j=Z_0=-1/2$.
The inverse Bernoulli transform gives the weights
\eqref{eg:ordinary-weights}. The centered expansion of the tail product
is
\[
 A_p(n)A_q(n)=u_0x_n^{2-p-q}+u_1x_n^{-p-q}
                         +u_2x_n^{-p-q-2}+O(x_n^{-p-q-4}).
\]
Since $\deg P_b\le2q-2$ and $p=q-4$, its nonnegative powers are
exactly $S_b(x_n)$; the remaining summand is $O(x_n^{-2})$.
The continued centered sums of $x_n^4,x_n^2,1$ are all zero,
so the ordinary sum equals the moment combination. Equivalently,
finite centered power summation gives these polynomial subtraction
terms zero constant coefficient at the integer cutoff. The polygamma
normalization follows from even $p,q$.

For \eqref{eg:top-coefficient}, set every row output in
\eqref{eg:coefficient-equation} except the last equal to zero.
A homogeneous solution before the final boundary equation has initial
vector proportional to $(Y_{b+1},-X_{b+1})$, since its coefficient
at $b+1$ must vanish. Its coefficient at $b$ is $-E_b$ times that
scalar, and its continued coefficient at $b+2$ is $E_{b+1}$ times
that scalar. Hence the last row gives
\[
        \frac{r_b}{Y_b^*}=\frac{E_b}{d_{b+2}E_{b+1}}.
\]
The row and column scalings above show that
$\lambda_{q-1}=4bE_b/(d_{b+2}E_{b+1})$. Multiplication by
$\binom{2q-1}{2q-2}=4b-1$ gives the stated positive coefficient.
\end{proof}

\subsection{A concrete evaluation at the new pair \texorpdfstring{$(8,12)$}{(8,12)}}

A centered difference keeps the integer coefficients readable. Put
\begin{align*}
 K(y)=y^{13}\bigl(&469839+1269227y^2-11819495y^4\\
                 &+56443363y^6+83672540y^8+12834702y^{10}\bigr),\\
 P(x)&=K(x+\tfrac12)-K(x-\tfrac12),\\
 S(x)&=\frac{2951981460x^4+51273021735x^2+118231913818}{770}.
\end{align*}
Then the exact ordinary identity is
\begin{equation}\label{eg:8-12-sum}
 \boxed{\displaystyle
 \sum_{n=0}^{\infty}\left\{
 \frac{P(n+\tfrac12)}{7!\,11!}
       \psi^{(7)}(n+1)\psi^{(11)}(n+1)-S(n+\tfrac12)\right\}
 =-\frac{790768509}{16}.}
\end{equation}
For example, the coefficient row
\[
 (469839,1269227,-11819495,56443363,83672540,12834702)
\]
multiplied by $M_{8,12}$ is exactly
$(790768509/8,0,0,0,0,0)$. The centered difference converts this row
of blocks into $P$, and the first three centered tail-product
coefficients give precisely $S$. Thus \eqref{eg:8-12-sum} follows
already from the preceding corollary.

There is also a finite telescoping proof, independent of the
continued-moment interpretation. Let $c_0,\ldots,c_5$ be the six
coefficients in the parentheses defining $K$, in increasing order,
and let
\[
 F_m(N)=\sum_{k=1}^N k^m
       =\frac{B_{m+1}(N+1)-B_{m+1}}{m+1}.
\]
Define the explicitly specified polynomials
\begin{align*}
 U_8(N)&=\sum_{j=0}^5c_jF_{2j+1}(N),&
 U_{12}(N)&=\sum_{j=0}^5c_jF_{2j+5}(N),\\
 V(N)&=-\frac{N}{6160}\bigl(
 2417202210N^4-5764920315N^3+42395905270N^2\notag\\
 &\hspace{36mm}-129829345800N+260578219219\bigr).
\end{align*}
The first two definitions imply
\[
 \Delta U_8(N)=K(N)N^{-12},\qquad
 \Delta U_{12}(N)=K(N)N^{-8},
\]
where $\Delta f(N)=f(N)-f(N-1)$. Direct polynomial expansion also
gives
\begin{equation}\label{eg:V-difference}
 \Delta V(N)=\frac{K(N)}{N^{20}}+
      \frac{U_8(N-1)}{N^8}+\frac{U_{12}(N-1)}{N^{12}}
                         -S(N-\tfrac12).
\end{equation}
In particular the negative powers cancel on the right. These are
finite rational-function identities, reproducible from the six
integer coefficients and the displayed Bernoulli polynomials.
Now put
\[
 \mathcal F(N)=K(N)A_8(N)A_{12}(N)
                 +U_8(N)A_8(N)+U_{12}(N)A_{12}(N)+V(N).
\]
Using $A_m(N-1)=A_m(N)+N^{-m}$ and the three difference identities,
all mixed terms cancel to give
\[
 \mathcal F(N)-\mathcal F(N-1)
 =P(N-\tfrac12)A_8(N-1)A_{12}(N-1)-S(N-\tfrac12).
\]
Since $\mathcal F(0)=0$, the partial sum of
\eqref{eg:8-12-sum} through $n=N-1$ is exactly $\mathcal F(N)$.
For the limiting step, use
\begin{align*}
 A_8(N)&=\frac1{7N^7}-\frac1{2N^8}+\frac2{3N^9}
                                      -\frac1{N^{11}}+O(N^{-13}),\\
 A_{12}(N)&=\frac1{11N^{11}}-\frac1{2N^{12}}+\frac1{N^{13}}
                                 -\frac{91}{30N^{15}}+O(N^{-17}).
\end{align*}
The three tail-containing terms of $\mathcal F$ have constant
coefficients, respectively,
\[
       -\frac{116518888243}{2310},\qquad
       -\frac{193499245}{336},\qquad
        \frac{350665911}{220}.
\]
Their sum is $-790768509/16$, $V$ has zero constant coefficient,
and the coefficients of $N,N^2,\ldots,N^5$ cancel exactly.
The remainders above contribute $O(N^{-1})$. This proves the
ordinary evaluation by a finite identity and an asymptotic limit.

\subsection{Scope and reproducible checks}

The exact script checks \eqref{eg:determinant} against the original
Bernoulli matrix for $3\le b\le18$, checks the independent rotation
matrix, and verifies \eqref{eg:top-coefficient} where a full inverse
is computed. It also checks 1,946 minors of order two for gaps
$c=2,\ldots,6$, and verifies the displayed ordinary evaluation by
an exact finite telescoper. These checks detect normalization errors;
the all-order proof is Theorem~\ref{eg:two-minors}, not the finite
arithmetic range.

The theorem about minors of order two cannot be promoted to a
fixed-sign assertion for larger boundary determinants. For example,
for $c=3$ the exact signs are
\[
 \det(R_{6+i,j})_{1\le i,j\le3}<0,\qquad
 \det(R_{7+i,j})_{1\le i,j\le3}>0.
\]
Both are nonzero, and both arrays satisfy all the strict negative
minor statements in Theorem~\ref{eg:two-minors}. Thus order-two
comparison alone does not prove the general conjecture. Combining
all currently established families, the still-unsettled range is
\[
             8\le p<q,\qquad q-p\ge6,
             \qquad p,q\ \hbox{even}.
\]
A useful next problem is to identify an exact, sign-flexible
nonvanishing invariant for the higher boundary determinants, or to
combine the fixed-gap recurrence with the complementary determinant
of size $p/2$ from the shifted-Hankel approach. The rank problem and
the arithmetic independence of the actual zeta coordinates remain
distinct questions.
